\documentclass[10pt,a4paper]{amsart}
\usepackage[margin=19mm]{geometry}
\usepackage{amsmath,amssymb,mathtools}
\usepackage[scr=rsfs,cal=euler]{mathalfa}
\usepackage{xfrac}
\usepackage[unicode,hidelinks,bookmarks=false]{hyperref}
\newcommand{\ie}{i.e.\ }
\newcommand{\cf}{cf.\ }
\newcommand{\resp}{respectively\ }
\newcommand{\Subsection}{\subsection}
\providecommand{\nobreakdash}{\nobreak\hbox{-}}
\setlength{\emergencystretch}{2em}
\multlinegap0pt
\usepackage{amssymb}
\usepackage{algorithm}\usepackage{algpseudocode}
\usepackage{float}
\newtheorem{proposition}{Proposition}
\newcommand{\PP}{\mathbb{P}}
\newcommand{\cov}{\mathrm{Cov}}
\title{Oracle-free Boltzmann Sampling for Powersets}
\author{Jean Peyen}
\date{Source: arXiv:2601.09508v3 (CC BY 4.0)}
\begin{document}
\maketitle

\noindent\emph{Source and licence.} Oracle-free Boltzmann Sampling for Powersets. Version arXiv:2601.09508v3 (CC BY 4.0), distributed by arXiv under the Creative Commons Attribution 4.0 International licence, http://creativecommons.org/licenses/by/4.0/, as recorded in the arXiv licence metadata for this exact version and shown on the abstract page https://arxiv.org/abs/2601.09508v3. Full licence notice: \texttt{licenses/arxiv-2601.09508-CC-BY-4.0.txt}.\par\smallskip

\noindent\textit{Editorial scope.} Counted unit: the article's proposition prop:ProbCov1 (source0.tex lines 103--109). The article's complete original proof of it is delivered with it (source0.tex lines 111--121, one \texttt{proof} environment of 11 lines). No mathematics is written, repaired or re-derived: the statement and the proof are the article's own lines, in the article's own notation, and no retained line is substituted or re-flowed.  No further source-local context is needed: every cross-reference of every retained line resolves inside the two delivered spans.  Nothing was omitted from any retained span, so the package carries no omission record.  No retained line contains a citation, so the wrapper carries no bibliography and no bibliography entry.  No retained line contains a figure environment, an image-inclusion command or drawing code, so no image is delivered and the package declares no asset dependency.  Editorial formatting only, in addition: an AMS \texttt{amsart} preamble that restates the article's own declarations and macros used by the retained lines, namely the environments \texttt{proposition} on separate counters, together with the standard packages those lines need.  The retained environments print in this document as Proposition 1 in order of appearance, whereas the article prints the same items with the numbering of its own full document.  The complete native source of the article as distributed in the arXiv e-print for this exact version is bundled unchanged as \texttt{sources/192683/source0.tex}.
\begin{proposition}\label{prop:ProbCov1}
\begin{align}
    \PP(\nu^{\prime}_\ell = 1) &= 1-G_M(1-f_\ell) \label{eq:Prob1}\\
    \cov(\nu^{\prime}_\ell, \nu^{\prime}_{\ell^\prime}) &= G_M(1-(f_\ell+f_{\ell^\prime}))-G_M(1-f_\ell)G_M(1-f_{\ell^\prime}) \label{eq:Cov1}
\end{align}
where $G_M$ is the probability generating function of $M$.
\end{proposition}

\begin{proof}
For \eqref{eq:Prob1},
\begin{equation*}
\PP(\nu^{\prime}_\ell = 1)=1-\PP(\nu_\ell=0)=1-\sum_{m\ge 0}\PP(M=m)(1-f_\ell)^m=1-G_M(1-f_\ell).
\end{equation*}
For \eqref{eq:Cov1}, using inclusion--exclusion,
\begin{equation*}
\PP(\nu^{\prime}_\ell \nu^{\prime}_{\ell^\prime}=1)=1-G_M(1-f_\ell)-G_M(1-f_{\ell'})+G_M(1-(f_\ell+f_{\ell'})),
\end{equation*}
so subtracting $\PP(\nu^{\prime}_\ell = 1)\PP(\nu^{\prime}_{\ell^\prime} = 1)$ gives \eqref{eq:Cov1}.
\end{proof}

\end{document}
